\chapter{Ergebnisse}

Zwischen einer Kapitelüberschrift und der ersten Unterüberschrift sollte immer ein einleitender Text stehen, der in das Kapitel einführt und dessen Inhalt und Struktur kurz beschreibt. Vermeiden Sie es, direkt von einer Überschrift zur nächsten zu springen, ohne den Lesenden Orientierung zu geben.

\section{Erster Abschnitt}
Im Kapitel \glqq{}Ergebnisse\grqq{} werden die Befunde der Untersuchung sachlich, nachvollziehbar und strukturiert dargestellt. Die Darstellung sollte sich an der Forschungsfrage oder der gewählten Methodik orientieren. Ziel ist es, die Lesenden in die Lage zu versetzen, die Ergebnisse zu verstehen und nachzuvollziehen, ohne dass bereits eine umfassende Interpretation oder Diskussion erfolgt.
Hinweise zur Darstellung:
\begin{itemize}
    \item \textbf{Strukturierung}: Gliedern Sie die Ergebnisse logisch, etwa nach Forschungsfragen, Hypothesen oder methodischen Schritten.
    \item \textbf{Objektivität}: Vermeiden Sie wertende Aussagen. Interpretationen gehören in die Diskussion oder das Fazit.
    \item \textbf{Visualisierung}: Nutzen Sie Tabellen, Abbildungen oder Diagramme, um komplexe Ergebnisse verständlich darzustellen. Jede Visualisierung muss im Text erläutert werden.
    \item \textbf{Nachvollziehbarkeit}: Beschreiben Sie die Ergebnisse so, dass die Lesenden die Schlussfolgerungen nachvollziehen können, ohne die Rohdaten einsehen zu müssen.
\end{itemize}
Falls die Arbeit eine klare Trennung zwischen Ergebnisdarstellung und Interpretation vorsieht, sollte diese hier eingehalten werden. In manchen Fachbereichen werden Ergebnisse und Diskussion getrennt, in anderen gemeinsam behandelt. Orientieren Sie sich an den fachspezifischen Konventionen und den Vorgaben Ihrer Betreuung.